% Part: computability
% Chapter: machines-computations
% Section: configuration

\documentclass[../../../include/open-logic-section]{subfiles}

\begin{document}

\olfileid{cmp}{tur}{con}
\olsection{مکمل حالتاں تے حسابی کارروائیاں}

\begin{explain}
پہلاں جفت مشین دیاں مکمل
حالتاں وچوں قدم بہ قدم
لنگھنا یاد کرو۔ پٹی،
پڑھن-لکھن والے سر تے
ٹیورنگ مشین دے پروگرام
توں بنیا تصوراتی نظام
دراصل ٹیورنگ مشین دی
حسابی کارروائی نوں سمجھݨ
دا اک بدیہی طریقہ اے۔
رسمی طور اُتے، دتے ہوئے
انپٹ لئی مشین دی حسابی
کارروائی نوں
\emph{مکمل حالتاں} دے سلسلے
وجوں تعریف کر سکدے آں۔
ہر مکمل حالت نشاناں
دے اک سلسلے (جیڑا
حسابی کارروائی دے اوس
ویلے پٹی دا متن دکھاندا
اے)، پڑھن-لکھن والے
سر دی جگہ دسن والا
اک عدد، تے مشین دی
اک حالت توں بندی اے۔ ایہناں دی
مدد نال اسی تعریف کر
سکدے آں کہ ٹیورنگ
مشین~$M$ دتے ہوئے
انپٹ اُتے کیہ حساب
کردی اے۔
\end{explain}

\begin{defn}[مکمل حالت]
ٹیورنگ مشین $M = \tuple{Q, \Sigma, q_0,
\delta}$ دی اک \emph{مکمل حالت}
تِن رکناں دا ترتیب وار
ٹپل $\tuple{C, m, q}$ اے،
جتھے:
\begin{enumerate}
\item $C \in \Sigma^*$، $\Sigma$
  دے نشاناں دا اک محدود
  سلسلہ اے،
\item $m \in \Nat$ اک عدد اے
  تے $< \len{C}$ اے، تے
\item $q \in Q$۔
\end{enumerate}
بدیہی طور اُتے سلسلہ~$C$
پٹی دا متن اے۔ ایہ سب توں کھبے خانے توں
آخری متعلقہ خانے تک دے سارے نشان نیں۔ آخری
متعلقہ خانہ اوہناں خانیاں وچوں سب توں سجّے پاسے اے
جیہڑے خالی نہیں یا جیہناں تک سر پہلے پہنچ چکیا اے۔ $m$~اوس
خانے دا نمبر اے جیہنوں
پڑھن-لکھن والا سر
ایس ویلے پڑھ رہیا اے؛
سب توں کھبے خانے دا
نمبر $0$ اے۔ $q$ مشین دی
موجودہ حالت اے۔
\end{defn}

\begin{explain}
ٹیورنگ مشین دا ممکنہ
انپٹ نشاناں دا اک سلسلہ
اے، عموماً اوہ سلسلہ
جیڑا کسے نہ کسے روپ
وچ اک عدد دا کوڈ اے۔ مشین دی
ابتدائی مکمل حالت
اوہ حالت اے جتھوں اسی
اُس انپٹ اُتے مشین
دا عمل شروع کردے آں:
پٹی اُتے پہلے پٹی دے
سِرے دا نشان تے اوہدے
فوراً مگروں سجے پاسے
دے خانیاں وچ انپٹ
لکھیا ہوندا اے؛
پڑھن-لکھن والا سر
انپٹ دے سب توں کھبے
خانے نوں پڑھ رہیا
ہوندا اے (یعنی کھبے
سِرے دے نشان توں
اگلے سجے خانے نوں)،
تے مشین مقرر کیتی
ابتدائی حالت~$q_0$
وچ ہوندی اے۔
\end{explain}

\begin{defn}[ابتدائی مکمل حالت]
$M$ دی انپٹ $I \in \Sigma^*$
لئی \emph{ابتدائی مکمل
حالت} ایہ اے:
\[
\tuple{\TMendtape \frown I, 1, q_0}.
\]
\end{defn}

\begin{explain}
نشان~$\frown$ دا مطلب
\emph{جوڑنا} (concatenation)
اے۔
% PNB-CMP-027: source B9 says left; B7 and the initial tuple put input right of the endmarker.
انپٹ دی لڑی پٹی دے
کھبے سِرے دے نشان
توں فوراً سجے پاسے
شروع ہوندی اے۔
\end{explain}

\begin{defn}
اسی آکھدے آں کہ مکمل
حالت $\tuple{C, m, q}$
\emph{اک قدم وچ مکمل حالت
$\tuple{C', m', q'}$ پیدا
کردی اے} (مشین~$M$ دے
مطابق)، اُدوں تے صرف
اُدوں جدوں:
\begin{enumerate}
\item $m$-ویں تھاں اُتے $C$ دا نشان
  $\sigma$ ہووے،
\item $M$ دے حکماں دا
  مجموعہ $\delta(q, \sigma) =
  \tuple{q', \sigma', D}$
  مقرر کردا ہووے،
\item $m$-ویں تھاں اُتے $C'$ دا
  نشان $\sigma'$ ہووے، تے
\item
\begin{enumerate}
\item جے $D = L$ ہووے تاں
  $m' = m - 1$، بشرطیکہ
  $m>0$ ہووے؛ نہیں
  تاں $m'=0$، یا
\item $D = R$ تے $m' = m + 1$، یا
\item $D = N$ تے $m' = m$،
\end{enumerate}
\item جے $m' = \len{C}$
  ہووے تاں $\len{C'} = \len{C} + 1$
  تے $m'$-ویں تھاں اُتے $C'$ دا
  نشان~$\TMblank$ ہووے۔
  نہیں تاں $\len{C'}=\len{C}$۔
\item ہر اوس $i$ لئی
  جدے لئی $i < \len{C}$
  تے $i \neq m$ ہووے،
  $C'(i) = C(i)$ ہووے،
\end{enumerate}
\end{defn}

\begin{defn}\ollabel{defn:run-output}
$M$ دا \emph{انپٹ~$I$
اُتے چلن} مکمل حالتاں
دا سلسلہ $C_i$ اے،
تے ایہ ساریاں مکمل حالتاں $M$ دیاں نیں۔
ایتھے $C_0$، $M$ دی
انپٹ~$I$ لئی ابتدائی مکمل
حالت اے، تے ہر $C_i$
اک قدم وچ $C_{i+1}$
پیدا کردی اے۔

اسی آکھدے آں کہ $M$
انپٹ $I$ اُتے $k$ قدماں
مگروں \emph{رُکدی اے}
جے $C_k =
\tuple{C, m, q}$ ہووے،
~$m$-ویں تھاں اُتے $C$ دا
نشان~$\sigma$ ہووے، تے
$\delta(q,
\sigma)$ غیر تعریف شدہ
ہووے۔ ایس صورت وچ
مشین~$M$ دے انپٹ~$I$ دا
\emph{آؤٹ پٹ}~$O$ اے؛
$O$ نشاناں دی ایداں دی
لڑی اے جیڑی~$\TMblank$
اُتے ختم نہیں ہوندی،
تے $C = \TMendtape \concat O \concat \TMblank^j$
کسے~$j \in \Nat$
لئی ہووے۔ ($\TMblank^j$،
$j$ عدد~$\TMblank$ نشاناں
دا سلسلہ اے۔)
\end{defn}

\begin{explain}
ایس تعریف دے مطابق
آؤٹ پٹ~$O$، جیہڑا مشین~$M$ دا اے،
ہمیشہ~$\TMblank$ توں
وکھرے نشان اُتے ختم
ہوندا اے؛ یا جے ویلے~$k$
پوری پٹی~$\TMblank$
نال بھری ہووے (سب توں
کھبے~$\TMendtape$
توں سوا)، تاں $O$
خالی لڑی اے۔
\end{explain}

\end{document}
